\documentclass[11pt]{article}
\usepackage[margin=0.9in]{geometry}
\usepackage{amsmath,amssymb,booktabs,enumitem,listings,xcolor,hyperref}
\definecolor{backcolour}{rgb}{0.95,0.95,0.92}
\definecolor{codeblue}{rgb}{0.05,0.15,0.55}
\definecolor{codegreen}{rgb}{0.0,0.45,0.0}
\lstset{
    backgroundcolor=\color{backcolour},
    basicstyle=\ttfamily\small,
    keywordstyle=\color{codeblue}\bfseries,
    commentstyle=\color{codegreen},
    stringstyle=\color{purple},
    breaklines=true,
    frame=single,
    showstringspaces=false,
    tabsize=4
}

\title{CPSC 250 \\ Philosophy of Complex Objects}
\author{Edward Brash}
\date{}
\begin{document}
\maketitle

\section{Motivation}

Before we begin writing our own classes, it is useful to think about why classes exist in the first place.

A program often needs to represent a ``thing'' from the real world. That thing may not be naturally described by a single integer, float, string, list, tuple, or dictionary. Instead, it may require several related pieces of information, possibly of different data types.

For example, imagine that we want to store information about a student.

\section{A Student as a Collection of Information}

A student might have the following data associated with them:

\begin{center}
\begin{tabular}{clc}
\toprule
Number & Information & Type \\
\midrule
1 & First name & string \\
2 & Last name & string \\
3 & Student ID & integer \\
4 & Address & list \\
5 & Grades & dictionary \\
\bottomrule
\end{tabular}
\end{center}

The address itself might be a list:

\begin{lstlisting}[language=Python]
address = [908, "Promenade Lane", "Williamsburg", "VA", 23185]
\end{lstlisting}

The grades might be a dictionary:

\begin{lstlisting}[language=Python]
grades = {
    "HW": 94.2,
    "Quiz": 75.8,
    "Tests": 88.4,
    "Final": 74.2
}
\end{lstlisting}

Already, this is no longer a simple primitive value. It is a structured collection of information.

\section{The List-Based Approach}

One possible approach is to store all the student information in a list:

\begin{lstlisting}[language=Python]
student_info = [first_name, last_name, student_id, address, grades]
\end{lstlisting}

Then we could access the homework grade by writing:

\begin{lstlisting}[language=Python]
print(student_info[4]["HW"])
\end{lstlisting}

This works, but it requires us to remember that index 4 contains the grades dictionary.

\section{Why This Works, But Is Not Ideal}

The list-based approach technically works.

However, it has several disadvantages:

\begin{itemize}
    \item The meaning of each index must be remembered.
    \item The code becomes difficult to read.
    \item The structure is fragile.
    \item If we change the order of the list, many parts of the program may break.
\end{itemize}

For example:

\begin{lstlisting}[language=Python]
student_info[4]["HW"]
\end{lstlisting}

requires us to remember that:

\begin{itemize}
    \item index 4 means ``grades'',
    \item and \verb|"HW"| is a key inside the grades dictionary.
\end{itemize}

This can quickly become confusing.

\section{A Complex Object}

The key idea is that \verb|student_info| is really a complex assembly of information.

It contains:

\begin{itemize}
    \item strings,
    \item integers,
    \item lists,
    \item dictionaries,
    \item and possibly even other objects.
\end{itemize}

It is natural to call this entire collection an \textbf{object}.

\section{Everything in Python Is Already an Object}

In Python, almost everything can be thought of as an object:

\begin{itemize}
    \item integers,
    \item floats,
    \item strings,
    \item booleans,
    \item lists,
    \item tuples,
    \item dictionaries,
    \item lists of lists,
    \item lists of dictionaries,
    \item dictionaries of lists,
    \item and so on.
\end{itemize}

Python already treats these things as objects internally.

The next step is to ask:

\begin{quote}
Can we define our own object types?
\end{quote}

The answer is yes.

\section{Main Point}

It would be useful to have a system that allows us to:

\begin{enumerate}
    \item define objects in our own way,
    \item decide what information belongs inside those objects,
    \item define how users interact with those objects,
    \item and protect the internal structure from accidental misuse.
\end{enumerate}

This is the motivation for classes and object-oriented programming.

\section{Classes and Object-Oriented Programming}

A class allows us to define a new kind of object.

Instead of forcing student data into a list, we could imagine defining a \verb|Student| class.

Then each student would be a \verb|Student| object, with meaningful internal variables such as first name, last name, ID number, address, and grades.

The advantage is that the program becomes much more readable and maintainable.

\section{Key Takeaways}

\begin{itemize}
    \item Real-world data often consists of many related pieces of information.
    \item Primitive types alone are not always enough.
    \item Lists and dictionaries can store complex information, but may become awkward.
    \item A complex assembly of information can be thought of as an object.
    \item Classes allow us to define our own object types.
    \item Object-oriented programming is about defining objects and controlling how users interact with them.
\end{itemize}

\end{document}
